\exercice{Intoxication au monoxyde de carbone (20 minutes)}\label{exo:intoxication_CO}
\points{8}
\medskip

Le monoxyde de carbone est un gaz inodore et invisible responsable de plus de la moitié des morts par intoxication dans le monde.
Issu d'une combustion incomplète, le monoxyde de carbone peut avoir des sources variées : chaudières au gaz défectueuse, poêles ou cheminées qui ne tirent pas assez, moteur thermique utilisé avec un mauvais carburant, départ de feu dans une maison, etc.

\begin{doc}{Détection du monoxyde de carbone}[\label{doc:detection_monoxyde}]
    Pour détecter le monoxyde de carbone, on peut acheter un détecteur spécialisé qui utilise une réaction d'oxydo-réduction avec les deux couples suivants :
    \begin{alignat*}{2}
      \dioxydeDeCarbone/\chemfig{CO} : \quad
      & 2\chemfig{CO} + 2\eau 
      && \reaction 2\dioxydeDeCarbone + 4\ionHydrogene + 4\electron \\
      %
      \dioxygene/\eau : \quad
      & 2\eau
      && \reaction \dioxygene + 4\ionHydrogene + 4\electron
    \end{alignat*}
  Les électrons échangés aux cours de la réaction d'oxydo-réduction entre le monoxyde de carbone et le dioxygène génèrent un courant électrique qui alimente une alarme sonore au delà d'un certain taux de monoxyde de carbone.

  Il est vivement recommandé d'avoir un détecteur de monoxyde de carbone chez soit, ils sont peu coûteux et peuvent éviter des tragédies.
\end{doc}

\question{
  Expliquer pourquoi on ne peut pas se baser sur notre vue ou notre odorat pour détecter le monoxyde de carbone.
  \points{1}
}{}

\question{
  Justifier que le monoxyde de carbone est un réducteur.
  \points{1}
}{}

\question{
  Écrire l'équation de réaction du monoxyde de carbone \chemfig{CO} avec le dioxygène \chemfig{O_2} en utilisant les deux demi-équations d'oxydo-réduction fournies.
  \points{3}
}{}


\begin{doc}{Effets physiologiques du \chemfig{CO}}[\label{doc:effets_physiologiques}]
  Quand nous respirons, le dioxygène est transporté des poumons aux tissus grâce à une protéine : l'hémoglobine notée Hb, ce qui permet la respiration cellulaire.
  Le monoxyde de carbone \chemfig{CO} se lie 250 fois plus facilement à l'hémoglobine que le dioxygène.
  Quand le monoxyde de carbone se lie à une ou plusieurs hémoglobine, on parle de carboxyhémoglobine notée COHb.

  \begin{listePoints}
    \item Si le taux de COHb est inférieur à \qty{10}{\percent}, il y a peu ou pas d'effets physiologiques observés.
    \item À \qty{15}{\percent} de COHb, une personne va commencer à avoir des maux de têtes.
    \item Entre \qty{20}{\percent} et \qty{30}{\percent}, les maux de têtes sont sévères et accompagnés de nausées, d'étourdissement, de confusion et de légères hallucinations visuelles.
    \item Entre \qty{30}{\percent} et \qty{50}{\percent} les symptômes neurologiques deviennent plus graves.
    \item Au delà de \qty{50}{\percent}, il y a perte de connaissance et un risque de tomber dans le coma, avec un risque d'insuffisance respiratoire.
  \end{listePoints}
  La mort arrive quand le taux de COHb dépasse les \qty{60}{\percent}.
\end{doc}

\question{
  Expliquer pourquoi le monoxyde de carbone est dangereux et pourquoi il faut rapidement évacuer un lieu si des symptômes apparaissent.
  \points{1,5}
}{}

\question{
  Pour traiter les intoxications, on fait respirer aux patient-es du dioxygène pur pressurisé pour faire redescendre rapidement le taux de COHb. Expliquer pourquoi il faut traiter rapidement les patient-es.
  \points{1,5}
}{}